% A卷 Q2 — 简答题
\begin{question}{0.1}
(20分) 简答题
\begin{enumerate}
\item 简述静定结构与超静定结构的主要区别（至少列出3点）。
\item 简述力法中利用对称性取半结构的方法（包括对称荷载和反对称荷载两种情况）。
\end{enumerate}
\end{question}

\begin{solution}
1.~区别：(1)静定结构全部反力和内力可由平衡方程唯一确定，超静定结构需补充变形协调条件；
(2)静定结构内力与材料性质和截面尺寸无关，超静定结构有关；
(3)温度变化和支座移动不引起静定结构内力，但引起超静定结构内力。

2.~对称结构对称荷载：反对称未知力为零，取半结构时对称轴处设为定向滑动支座。
对称结构反对称荷载：对称未知力为零，取半结构时对称轴处设为活动铰支座或链杆。
\end{solution}
